\subsection{\problemblock{*Print} Data Block}\label{sec:problem-block-print}

The \problemblock{*print} data block can appear either within a trajectory or within
a problem (outside of a trajectory definition). The purpose of the
\problemblock{*print} data block is the same in either case: to write trajectory
information to the standard output file. The standard output file has the same
filename as the problem file, but it is of type \texttt{.out} instead of
\texttt{.prb}. It has page breaks and headings formatted for a standard printer with
132 characters per line and 60 lines per page.

The difference between \problemblock{*print} data blocks located within a
trajectory definition and those located outside of a trajectory is that those located
within a trajectory assume the printed information is for that trajectory. The
variables printed are assumed to be associated with the trajectory containing the
\problemblock{*print} data block.

\problemblock{*print} data blocks that are not within a trajectory definition do not
make this assumption, so trajectory numbers enclosed in brackets are required in
addition to the variable names. This provides the capability to print variables from
several trajectories on the same page.

For example,

\begin{lstlisting}[style=taosmanual]
*print
  time[1] alt[1] alt[2] range[1] range[2]
\end{lstlisting}

prints the time from trajectory number 1, the altitude from trajectories 1 and 2, and
the range from trajectories 1 and 2 as shown in
\cref{fig:problem-printout-example}. This makes it easy to compare values from
different trajectories. Data from different trajectories are related by mission time so
\statevar{alt[1]} and \statevar{alt[2]} refer to the altitudes of trajectories 1 and 2
at the same instant in time.

\begin{center}
  \begin{minipage}{0.92\linewidth}
    \lstinputlisting[style=taosmanual,basicstyle=\ttfamily\scriptsize]{examples/chapter04/problem-printout-example.txt}
    \captionof{figure}{Trajectory printout example.}
    \label{fig:problem-printout-example}
  \end{minipage}
\end{center}

The printout begins with a page heading, a title, and column headings that are
repeated on every page. This is followed by the trajectory variables. Any of the
standard output variables listed in the Appendix can be printed along with any
user-defined variables. The output format of each variable is controlled with the
\problemblock{*units/fmt} data block (\cref{sec:block-units-fmt}).

A table of trajectory data is written, with a maximum of 10 output variables, for each
\problemblock{*print} data block. Any number of \problemblock{*print} data
blocks can be included within a problem.

Radar and relative-vehicle output variables, those starting with \texttt{rad} or
\texttt{rel}, must contain an additional subscript enclosed in square brackets. The
first subscript gives the radar station number or the trajectory number for the
relative-vehicle calculations, and the second subscript gives the vehicle trajectory
number. For example, the variable \statevar{radrng[2][1]} is used to print the radar
range of trajectory number 1 from radar station number 2, and the variable
\statevar{relvel[4][2]} is used to print the closing velocity of trajectory number 4
relative to trajectory number 2.
